%=============================================================================
% Chapter 13: API Documentation
%=============================================================================
\chapter{API Documentation}

This chapter documents all implemented API endpoints in FloodEye. Each endpoint is a Next.js Route Handler deployed as a Vercel serverless function under \texttt{app/api/}.

\section{Base URL}

\begin{itemize}[leftmargin=1.5em]
  \item \textbf{Development:} \texttt{http://localhost:3000}
  \item \textbf{Production:} \texttt{https://\{project\}.vercel.app}
\end{itemize}

\section{Common Response Headers}

All endpoints return:
\begin{itemize}[leftmargin=1.5em]
  \item \texttt{Content-Type: application/json}
  \item ISR cache headers set by the \texttt{revalidate} export.
\end{itemize}

\section{Endpoint: GET /api/sensors}

\subsection{Purpose}
Returns the single most recent sensor reading from the ESP32 via ThingSpeak, along with device status information.

\subsection{Cache}
\texttt{export const revalidate = 15;} --- Responses are cached by Vercel CDN for 15 seconds.

\subsection{Request}
\begin{lstlisting}[language={},caption={GET /api/sensors request}]
GET /api/sensors
Headers: (none required)
Body: (none)
\end{lstlisting}

\subsection{Response (200 OK)}
\begin{lstlisting}[language={},caption={GET /api/sensors success response}]
{
  "data": {
    "temperature": 24.5,
    "humidity": 67.0,
    "pressure": 1012.3,
    "altitude": 118.4,
    "distance": 82.0,
    "timestamp": "2026-07-10T10:30:00Z"
  },
  "deviceStatus": {
    "esp32Online": true,
    "thingspeakConnected": true,
    "lastUpload": "2026-07-10T10:30:00Z",
    "rssi": 0
  }
}
\end{lstlisting}

\subsection{Response (503 Service Unavailable)}
Returned when ThingSpeak credentials are missing or the ThingSpeak request fails.
\begin{lstlisting}[language={},caption={GET /api/sensors error response}]
{
  "error": "Failed to fetch data from ThingSpeak. Check your credentials in .env.local."
}
\end{lstlisting}

\subsection{Status Codes}
\begin{tabular}{|l|l|}
\hline
\rowcolor{FloodLight}
\textbf{Status} & \textbf{Meaning} \\
\hline
200 & Success \\
503 & ThingSpeak unavailable or credentials missing \\
\hline
\end{tabular}

\section{Endpoint: GET /api/history}

\subsection{Purpose}
Returns an array of historical sensor readings from ThingSpeak for a specified time range.

\subsection{Cache}
\texttt{export const revalidate = 60;} --- Responses cached for 60 seconds.

\subsection{Query Parameters}
\begin{tabular}{|l|l|l|}
\hline
\rowcolor{FloodLight}
\textbf{Parameter} & \textbf{Type} & \textbf{Values} \\
\hline
range & string & \texttt{1h}, \texttt{6h}, \texttt{24h}, \texttt{7d}, \texttt{30d} (default: \texttt{1h}) \\
\hline
\end{tabular}

\subsection{Request}
\begin{lstlisting}[language={},caption={GET /api/history request}]
GET /api/history?range=6h
\end{lstlisting}

\subsection{Response (200 OK)}
\begin{lstlisting}[language={},caption={GET /api/history success response}]
[
  {
    "temperature": 24.1,
    "humidity": 65.0,
    "pressure": 1013.1,
    "altitude": 118.0,
    "distance": 84.0,
    "timestamp": "2026-07-10T04:30:00Z"
  },
  ...
]
\end{lstlisting}

\subsection{Range to Query Mapping}
\begin{table}[H]
\centering
\caption{History Range Query Strategy}
\begin{tabular}{|l|l|l|}
\hline
\rowcolor{FloodLight}
\textbf{Range} & \textbf{Strategy} & \textbf{ThingSpeak Parameter} \\
\hline
1h & Count-based & \texttt{results=240} \\
6h & Count-based & \texttt{results=1440} \\
24h & Count-based & \texttt{results=5760} \\
7d & Date range & \texttt{start=...\&end=...} \\
30d & Date range & \texttt{start=...\&end=...} \\
\hline
\end{tabular}
\end{table}

\section{Endpoint: GET /api/analytics}

\subsection{Purpose}
Returns per-sensor descriptive statistics (average, min, max, trend) computed over the specified time range.

\subsection{Cache}
\texttt{export const revalidate = 60;}

\subsection{Query Parameters}
Same \texttt{range} parameter as \texttt{/api/history}. Default: \texttt{24h}.

\subsection{Request}
\begin{lstlisting}[language={},caption={GET /api/analytics request}]
GET /api/analytics?range=24h
\end{lstlisting}

\subsection{Response (200 OK)}
\begin{lstlisting}[language={},caption={GET /api/analytics success response}]
{
  "range": "24h",
  "count": 5412,
  "analytics": {
    "temperature": { "avg": 24.3, "min": 21.1, "max": 28.7, "trend": "up" },
    "humidity":    { "avg": 68.2, "min": 55.0, "max": 82.1, "trend": "stable" },
    "pressure":    { "avg": 1011.5, "min": 1008.0, "max": 1014.2, "trend": "down" },
    "altitude":    { "avg": 118.4, "min": 117.0, "max": 120.1, "trend": "stable" },
    "distance":    { "avg": 79.3,  "min": 35.0, "max": 110.0, "trend": "up" }
  }
}
\end{lstlisting}

\subsection{Trend Computation}
The trend is computed by comparing the average of the last 10\% of readings against the first 10\%:
\begin{itemize}[leftmargin=1.5em]
  \item \textbf{up}: last average $>$ first average $+ 0.5$
  \item \textbf{down}: last average $<$ first average $- 0.5$
  \item \textbf{stable}: otherwise
\end{itemize}

\section{Endpoint: GET /api/environment-score}

\subsection{Purpose}
Returns the computed Environmental Comfort Score based on the latest temperature and humidity readings.

\subsection{Request}
\begin{lstlisting}[language={},caption={GET /api/environment-score request}]
GET /api/environment-score
\end{lstlisting}

\subsection{Response (200 OK)}
\begin{lstlisting}[language={},caption={GET /api/environment-score response}]
{
  "score": 78,
  "status": "Good"
}
\end{lstlisting}

Score ranges: 0--40 (Poor), 41--60 (Moderate), 61--80 (Good), 81--100 (Excellent).

\section{Endpoint: POST /api/flood-predict}

\subsection{Purpose}
Structural placeholder for future server-side TensorFlow.js Node.js inference. In the current implementation, actual inference runs client-side via \texttt{lib/flood-model.ts}.

\subsection{Request Body}
\begin{lstlisting}[language={},caption={POST /api/flood-predict request body}]
{
  "history": [
    { "temperature": 24.5, "humidity": 67.0, "pressure": 1012.3,
      "altitude": 118.0, "distance": 82.0,
      "timestamp": "2026-07-10T09:30:00Z" },
    ...  // minimum 48 entries required
  ]
}
\end{lstlisting}

\subsection{Response (501 Not Implemented)}
\begin{lstlisting}[language={},caption={POST /api/flood-predict current response}]
{
  "message": "Server-side inference is currently disabled. Please run predictions client-side using lib/flood-model.ts",
  "status": "pending_implementation"
}
\end{lstlisting}

\subsection{Error Responses}
\begin{tabular}{|l|l|}
\hline
\rowcolor{FloodLight}
\textbf{Status} & \textbf{Condition} \\
\hline
400 & Invalid payload (missing or non-array \texttt{history} field) \\
400 & Insufficient data (\texttt{history.length < 48}) \\
501 & Server-side inference not yet implemented \\
500 & Internal server error \\
\hline
\end{tabular}

\section{API Authentication}

The current implementation does not implement API-level authentication. Environment variables (\texttt{THINGSPEAK\_CHANNEL\_ID}, \texttt{THINGSPEAK\_READ\_API\_KEY}) are kept server-side only and are never exposed to the client. The ThingSpeak Read API Key provides read-only access to the channel and does not expose write credentials.

\section{Environment Variables}

\begin{table}[H]
\centering
\caption{Required Environment Variables}
\begin{tabularx}{\linewidth}{|l|l|X|}
\hline
\rowcolor{FloodLight}
\textbf{Variable} & \textbf{Scope} & \textbf{Description} \\
\hline
\texttt{THINGSPEAK\_CHANNEL\_ID} & Server-only & ThingSpeak channel numeric ID \\
\texttt{THINGSPEAK\_READ\_API\_KEY} & Server-only & ThingSpeak read-only API key \\
\texttt{NEXT\_PUBLIC\_MAPTILER\_KEY} & Public (client) & MapTiler API key for map tile service \\
\hline
\end{tabularx}
\end{table}
